\documentclass[reqno]{amsart} 
\usepackage{amssymb,latexsym,amsmath,amscd,graphicx,setspace,amsthm,verbatim}
\usepackage[margin = 3 cm]{geometry}


\theoremstyle{plain}
\newtheorem{theorem}{Theorem}[section]
\newtheorem{proposition}{Proposition}
\newtheorem{corollary}{Corollary}
\newtheorem{lemma}{Lemma}
\newtheorem{conjecture}{Conjecture}
\newtheorem{question}{Question}
\newtheorem{problem}{Problem}
      
\theoremstyle{definition}
\newtheorem{definition}{Definition}

\newenvironment{solution}{\paragraph{\emph{Solution}.}}{\hfill$\square$}


\newenvironment{solution1}{\paragraph{\emph{Solution $1$}.}}{\hfill$\square$}
\newenvironment{solution2}{\paragraph{\emph{Solution $2$}.}}{\hfill$\square$}
\newenvironment{solution3}{\paragraph{\emph{Solution $3$}.}}{\hfill$\square$}

\begin{document} 

\title[Homework 7]{Homework 7 \\ First due date: October 29.}

\date{\today} 
\maketitle 


\begin{problem}
Consider the smooth function $f:\mathbb{R}^{3} \rightarrow \mathbb{R}$ given by $(x,y,z) \mapsto f(x,y,z) = z^{2} - x^{2} - y^{2} -1$.  
\begin{enumerate}
\item Is $0 \in \mathbb{R}$ a regular value for $f$?  
\item Is 
$$S = f^{-1}(0) =  \{(x,y,z) \in \mathbb{R}^{3} : z^{2} - x^{2} - y^{2} = 1\} $$
a regular smooth manifold in $\mathbb{R}^{3}$?  If yes, what is its dimension?  
\item Can you draw (sketch) this smooth manifold?
\end{enumerate}
\end{problem}
\begin{solution1}

\end{solution1}

\begin{problem}
Although the theorem on regular level sets from class (a consequence of the implicit function theorem) allows us to bypass local parametrizations, could you find two (Monge) local parametrizations that would cover the regular smooth manifold from the previous problem??
\end{problem}
\begin{solution1}

\end{solution1}

\begin{problem}
Go back to the third way of defining a structure of a regular 1-dimensional smooth manifold on $S^{1}$ via the two local parametrizations $\varphi_{N}: \mathbb{R} \rightarrow V_{N} \cap S^{1}$ and $\varphi_{S}: \mathbb{R} \rightarrow V_{S} \cap S^{1}$.  (See Problem $1$ on Homework $5$.)  
\begin{enumerate}
\item Go ahead and find the two transition maps $\varphi_{S}^{-1} \circ \varphi_{N}$ and $\varphi_{N}^{-1} \circ \varphi_{S}$ paying particular attention to their domain and codomain.  
\item Are those two transition maps smooth?
\end{enumerate}
\end{problem}
\begin{solution1}

\end{solution1}


\begin{problem}
Go back to Problem $2$ of Homework $5$, pick your favorite two local parametrizations out of six, and find out the corresponding transition maps.  Are they smooth?
\end{problem}
\begin{solution1}

\end{solution1}



\begin{problem}
Consider the smooth function $f:\mathbb{R} \rightarrow \mathbb{R}$ given by $x \mapsto f(x) = x^{2}$.  We showed in class that its graph $\Gamma_{f}$ is a regular $1$-dimensional smooth manifold in $\mathbb{R}^{2}$ that can be covered by one (Monge) local parametrization $\varphi:\mathbb{R} \rightarrow \mathbb{R}^{2} \cap \Gamma_{f}$ given by
$$x \mapsto \varphi(x) = (x,f(x)). $$
Let $p = (2,4) \in \Gamma_{f}$.  
\begin{enumerate}
\item Find a basis for $T_{p}(\Gamma_{f})$.  
\item What is the dimension of the vector space $T_{p}(\Gamma_{f})$?
\end{enumerate}
\end{problem}
\begin{solution1}

\end{solution1}


\begin{problem}
Go back to Problem $2$ of Homework $5$ and consider $p = \left(\frac{1}{2},\frac{1}{2},\frac{1}{\sqrt{2}}\right) \in S^{2}$.
\begin{enumerate}
\item Pick one local parametrization that covers the point $p$.  Using that local parametrization, find a basis for $T_{p}(S^{2})$.
\item Pick a different local parametrization that covers the point $p$.  Using that new local parametrization, find a basis for $T_{p}(S^{2})$.
\item In class, we showed that the tangent space does {\bfseries not} depend on the choice of the local parametrization in order to calculate it.  Put your Math 235 hat on, and show carefully that the two bases you obtained above span the {\bfseries same} subspace of $\mathbb{V}^{3}$.
\end{enumerate}
\end{problem}
\begin{solution1}

\end{solution1}















\end{document} 


